\section{Question 2: deciding connectivity from $A$}\label{sec:connectivity}

\subsection{Connected graphs and components}

\plain In some countries you can drive from any town to any other town.
In others, some town is cut off from the rest.

\begin{definition}[connected; textbook Definition 2.1.1]\label{def:connected}
A graph $G$ is \emph{connected} if for any two vertices $u$ and $v$ there is
a path between $u$ and $v$ \cite[Definition 2.1.1, p.~48]{petrosyan2015grafneri}.
The board states the same definition: a graph is called connected if for any
$n_i,n_j$ there exists a path between the nodes \cite{whiteboard1lecture}.
\end{definition}

\begin{remark}
The textbook's word ``path'' in Definition 2.1.1 is its trail, while the
board and the course use the path of \cref{def:path}. \Cref{cor:readings}
shows that this does not change which graphs are connected: a walk, a trail
and a path between two vertices exist for the same pairs. A graph with one
vertex is connected, since there are no two different vertices to join.
\end{remark}

\ex The board map of \cref{fig:three}(a) is connected: the paths $n_1,n_2$ and
$n_2,n_3$ and $n_1,n_2,n_3$ join all pairs of distinct vertices, and every
vertex is joined to itself by a path of length $0$.
\nonex The map of \cref{fig:three}(c) is not connected: $n_3$ has no
neighbour, so no walk and hence no path joins $n_3$ to $n_1$
(\cref{cor:readings}).

\medskip
\plain Towns that can reach each other form an island. Towns on different
islands cannot reach each other. A country with one island is one where every town can be reached from every other.

\begin{definition}[connectivity component; textbook, p.~49]\label{def:component}
The relation ``there is a path between $u$ and $v$'' is an equivalence
relation on $V$ (reflexive, symmetric and transitive). Its classes induce
subgraphs of $G$, called the \emph{connectivity components} of $G$
\cite[p.~49]{petrosyan2015grafneri}. An equivalence class is a group of vertices that can all reach one another.
The relation is reflexive because a vertex reaches itself by a path of length
$0$, symmetric because a path can be read backwards, and transitive because a
walk from $u$ to $v$ followed by a walk from $v$ to $w$ is a walk from $u$ to
$w$, which contains a path by \cref{cor:readings}. The subgraph induced by a
set $X$ of vertices has the vertices in $X$ and the edges of $G$ with both ends
in $X$.
\end{definition}

\begin{remark}\label{rem:onecomp}
A graph is connected if and only if it has exactly one component
(\cite[Remark 2.1.1, p.~49]{petrosyan2015grafneri}; also
\cite[Definition 2.9.12, p.~42]{grinberg2025introduction}). Two vertices lie in
the same component exactly when a path, and by \cref{cor:readings} a walk,
joins them.
\end{remark}

\ex The map of \cref{fig:three}(c) has two components: the one with vertices
$n_1,n_2$ and the road $n_1n_2$, and the one with the single vertex $n_3$.
\nonex $\{n_1,n_3\}$ is not the vertex set of a component of the map of
\cref{fig:three}(c), since no path joins $n_1$ and $n_3$.

\subsection{The test}

\plain To see whether one town can reach another, look at trips of one road,
of two roads, and so on. Add up the tables for all these trips. A zero in a
place means no such trip joins the two towns, whatever the number of roads.
We only need to try a few numbers of roads. A trip that visits no town twice
has at most one road less than there are towns.

\begin{theorem}[connectivity test]\label{thm:test}
Let $G$ be a simple graph with $N\ge2$ vertices and adjacency matrix $A$, and
put $S=A+A^2+\dots+A^{N-1}$. Then $G$ is connected if and only if
$S_{ij}>0$ for every $i\ne j$.
\end{theorem}

The textbook has no matrix test (\cref{sec:sources}). Paramadevan and
Sotheeswaran state this test as their Theorem~2
\cite[p.~17]{paramadevan2021properties}. The proof below is ours and uses
\cref{thm:count}, \cref{lem:shorten} and \cref{prop:distance}(c).

\begin{proof}
\emph{Connected implies $S_{ij}>0$.} Let $G$ be connected and let $i\ne j$.
By \cref{def:connected} there is a path between $n_i$ and $n_j$ (by
\cref{cor:readings}, whichever of the readings of ``path'' is used); let $k$
be its length. Since $i\ne j$, $k\ge1$, and by \cref{prop:distance}(c)
$k\le N-1$. A path is a walk, so by \cref{thm:count} $(A^k)_{ij}\ge1$. Every
power of $A$ has non-negative entries (\cref{def:power}), so
$S_{ij}\ge(A^k)_{ij}\ge1$.

\emph{$S_{ij}>0$ for all $i\ne j$ implies connected.} Fix $i\ne j$. The number
$S_{ij}$ is the sum of the non-negative integers $(A^k)_{ij}$ for
$k=1,\dots,N-1$, so one of them is positive: some $k$ with $1\le k\le N-1$
has $(A^k)_{ij}\ge1$. By \cref{thm:count} there is a walk of length $k$ from
$n_i$ to $n_j$, and by \cref{lem:shorten} a path between $n_i$ and $n_j$. A
vertex is joined to itself by the path of length $0$. So any two vertices are
joined by a path, and $G$ is connected.
\end{proof}

\begin{remark}\label{rem:gap}
The step from a positive entry to a path is where care is needed. The entry
counts walks (\cref{prop:notpaths}), and a walk need not be a path.
\Cref{lem:shorten} supplies the step, and \cref{prop:distance}(c) shows that
powers up to $N-1$ are enough. The published proof of the same test omits this
step (\cref{sec:discussion}). For $N=1$ the sum $S$ is empty and the
single-vertex graph is connected, so the test is stated for $N\ge2$.
\end{remark}

\plain There is a neater way to write the same test. Add a stay-at-home
option to the road table. Then multiply that table by itself a number of
times, one less than the number of towns. The test then reads: no zero
anywhere.

\begin{corollary}[the form with $(I+A)^{N-1}$]\label{cor:IA}
A simple graph $G$ with $N\ge1$ vertices is connected if and only if the
matrix $(I+A)^{N-1}$ has no zero entry.
\end{corollary}

The idea is this. The binomial theorem holds for matrices that commute, and $I$
commutes with every matrix, so $(I+A)^{N-1}$ expands into $\sum_{k=0}^{N-1}\binom{N-1}{k}A^k$, a sum with positive
coefficients of the same powers as in $S$ plus $I$. The exponent $N-1$ is the same
bound as in $S$: it is the longest length a path can have. The diagonal is positive
because of $I$, and an off-diagonal entry is positive exactly when the same
entry of $S$ is. The full proof is in \cref{app:derivations}. No source that we
read states the test in this form with a proof; the result is derived in this
note.

\subsection{Worked examples}

\plain First the board country, with three towns in a row. We list the road
table, the table for trips of two roads, and their sum. Then we look for zeros.

\begin{example}[a connected map]\label{ex:connected}
The board map has $N=3$, so $S=A+A^2$ and the other form is $(I+A)^2$. The
matrices are
\[
A=\begin{pmatrix}0&1&0\\1&0&1\\0&1&0\end{pmatrix},\quad
A^2=\begin{pmatrix}1&0&1\\0&2&0\\1&0&1\end{pmatrix},\quad
S=\begin{pmatrix}1&1&1\\1&2&1\\1&1&1\end{pmatrix},
\]
\[
I+A=\begin{pmatrix}1&1&0\\1&1&1\\0&1&1\end{pmatrix},\qquad
(I+A)^2=\begin{pmatrix}2&2&1\\2&3&2\\1&2&2\end{pmatrix}.
\]
Every off-diagonal entry of $S$ is positive and $(I+A)^2$ has no zero entry,
so by \cref{thm:test} and \cref{cor:IA} the map is connected, as it should be.
The entry $S_{13}=1$ comes from $(A^2)_{13}=1$, the walk $n_1,n_2,n_3$. The
entry $(A^2)_{12}=0$ although $n_1$ and $n_2$ are adjacent; the entry $A_{12}=1$
makes up for it in $S$. This is why the powers are summed and a single power
is not enough. As a check of \cref{cor:IA}, $(I+A)^2=I+2A+A^2$ here.
\end{example}

\plain Now a country where one town has no road. We do the same
calculation and find zeros.

\begin{example}[a map that is not connected]\label{ex:disconnected}
Take the map of \cref{fig:three}(c), with the single road $n_1n_2$ and $n_3$
alone. Then
\[
A=\begin{pmatrix}0&1&0\\1&0&0\\0&0&0\end{pmatrix},\quad
A^2=\begin{pmatrix}1&0&0\\0&1&0\\0&0&0\end{pmatrix},\quad
S=\begin{pmatrix}1&1&0\\1&1&0\\0&0&0\end{pmatrix},
\]
\[
I+A=\begin{pmatrix}1&1&0\\1&1&0\\0&0&1\end{pmatrix},\qquad
(I+A)^2=\begin{pmatrix}2&2&0\\2&2&0\\0&0&1\end{pmatrix}.
\]
The matrix $S$ has zeros at $(1,3)$, $(2,3)$ and their mirror images, and
$(I+A)^2$ has zeros in the same places. By \cref{thm:test} and \cref{cor:IA} the
map is not connected. By \cref{cor:components} below the components are
$\{n_1,n_2\}$ and $\{n_3\}$.
\end{example}

\plain For the drawing on the board, read with a third road, every pair of
towns is joined by a road. We expect no zeros.

\begin{example}[the drawing read as three roads]\label{ex:triangle}
For the map of \cref{fig:three}(b) every entry of $S$ is $2$ and every entry
of $(I+A)^2$ is $3$. The map is connected, which is also visible directly.
\end{example}

\subsection{Components, distance and sharpness}

\plain The table of sums also tells which towns share an island. A zero in a
place means the two towns are on different islands.

\begin{corollary}[components from $S$]\label{cor:components}
For $i\ne j$, the vertices $n_i$ and $n_j$ lie in the same component of $G$
if and only if $S_{ij}>0$. So the pairs with $S_{ij}=0$ are exactly the pairs
in different components.
\end{corollary}

\begin{proof}
If $n_i$ and $n_j$ are in the same component, a path joins them
(\cref{rem:onecomp}), and the first half of the proof of \cref{thm:test}, applied
to this pair alone, gives $S_{ij}>0$. If $S_{ij}>0$, the second half of that
proof gives a path between $n_i$ and $n_j$, so they are in the same class of
the relation of \cref{def:component}.
\end{proof}

\begin{remark}[distance]\label{rem:distance}
For distinct vertices in the same component, $d(n_i,n_j)$ is the least $k\ge1$
with $(A^k)_{ij}>0$. Indeed, a walk of length $k$ exists exactly when the
entry is positive (\cref{thm:count}), and the shortest walk has the length of
the shortest path (\cref{prop:distance}(a)). On the board map,
$A_{13}=0$ and $(A^2)_{13}=1$, so $d(n_1,n_3)=2$, in agreement with
\cref{ex:boardpaths}.
\end{remark}

\plain Take towns in a row, each joined only to the next. To get from the first
town to the last, you must pass every road in between. So no shorter trip
exists. If we only looked at shorter trips, we would wrongly say the first and
last towns are cut off.

\begin{proposition}[the limit $N-1$ cannot be lowered]\label{prop:sharp}
Let $N\ge2$ and let $G$ be the path graph with vertices $n_1,\dots,n_N$ and edges
$n_in_{i+1}$ for $1\le i<N$. Then $(A^k)_{1N}=0$ for $k<N-1$ and
$(A^{N-1})_{1N}=1$. Hence a sum $A+\dots+A^m$ with $m<N-1$ has a zero at
position $(1,N)$, although $G$ is connected.
\end{proposition}

\begin{proof}
Each edge of $G$ joins two vertices whose indices differ by exactly $1$, so
each step of a walk changes the index by $+1$ or $-1$. A walk from $n_1$ to
$n_N$ must change the index from $1$ to $N$, a total change of $N-1$, so it
has at least $N-1$ steps. By \cref{thm:count} this gives $(A^k)_{1N}=0$ for
$k<N-1$. A walk of length $N-1$ must use $+1$ at every step, so the only one is
$n_1,n_2,\dots,n_N$, and $(A^{N-1})_{1N}=1$.
\end{proof}

A computer run on the path graphs with $N=3,4,5,6$ gave the first positive power
$A^{N-1}$ each time, as the proposition says \cite{m3connectivityscript}.

\subsection{Breadth-first search, for comparison}

\plain Stand in town 1. Visit all its neighbour towns. Then visit the
neighbours of those. Note each town once. If every town was visited, then
every town can be reached from town 1.

This section is a comparison and a cross-check; the note does not rely on it.
\emph{Breadth-first search} starts with the set $\{n_1\}$ and repeatedly adds
every $n_j$ with $A_{mj}=1$ for some $n_m$ already in the set. When nothing can
be added, $G$ is connected exactly when the set contains all $N$ vertices. The
vertices added are those joined to $n_1$ by a path, so by \cref{rem:onecomp} the
set is the component of $n_1$. Grinberg states that in a breadth-first search
tree the tree path from the root is a shortest path
\cite[p.~176]{grinberg2025introduction}; using the search to decide
connectivity is our reading, not a statement of that source. By our own count,
forming $S$ takes $N-2$ matrix products: $A^2,\dots,A^{N-1}$, each from the
previous one by one product. The search reads the row of each visited vertex at
most once. These are different units, so we make no claim about running time.
On the maps of \cref{ex:connected} and \cref{ex:disconnected} the search visits
$n_1,n_2,n_3$ and $n_1,n_2$ respectively \cite{m3connectivityscript}.

\subsection{A computer check}

This is an experimental finding, not a proof. A program listed every labelled
simple graph with $N=1,2,\dots,6$ vertices (1, 2, 8, 64, 1024 and 32768 graphs,
of which 1, 1, 4, 38, 728 and 26704 are connected). For each graph it decided
connectivity by the sum test, by the $(I+A)^{N-1}$ test and by breadth-first
search from $n_1$ (the previous subsection), and the three answers always
agreed \cite{m3connectivityscript}. The proofs above do not rest on this check.
